\section{Causally autonomous column groups}\label{sec:causal-groups}

The source-cover results restrict the geometry of complete dependence.
This section gives a complementary restriction on the functions computed
along that geometry. If every local update obeys a capacity-conserving
decomposition into fixed narrow groups of consecutive columns after one
helper is fixed, then a
linear-length computation must disclose almost the entire input through
that helper.  Neither the return hypotheses nor final column purity
alone supplies this causal decomposition.

\subsection{The causal hypothesis and the routing convention}

Here $b\ge1$ is any integer, $W=db$, and $L=d^2b$.  Consider a directed
acyclic network with nodes $0,\ldots,n-1$, all edges directed from smaller
to larger nodes, and all memory edges $(v,v+1)$.  Edge $e$ has a fixed
nonnegative capacity $c_e$, measured in bits.  Each input row $i$ has a
specified set $A_i$ of attachment nodes; at a node in $A_i$ the entire
raw row is visible.  Each raw bit $(i,j,t)\in[d]^2\times[b]$ has a
specified output node $o(i,j,t)$.  These nodes respect transposed output
order: the bits of column $j$ are output before those of column $j+1$,
although several columns can share a node.  These are the row-source
and output conventions of the block-network model
in~\cite[section ``Codes on time--space diagrams'']{FirstPaper2026}.

Let $X\in\{0,1\}^{L}$ have any probability law supported on correctly
transposed inputs, and let $Z$ be a deterministic function of $X$.
Partition the columns into nonempty consecutive intervals
$J_1,\ldots,J_h$, in increasing order.  Write $g_a=|J_a|$,
$g=\max_a g_a$, and $X_a=X|_{[d]\times J_a\times[b]}$.
All entropies and logarithms in this section use base~$2$.

\Needspace{17\baselineskip}
\begin{definition}[Causal group typing]\label{def:causal-group-typing}
\formal{SourceCovers/CausalGroups.lean}{SourceCovers.CausalGroups.ColumnPartition,
SourceCovers.CausalGroups.Network, SourceCovers.CausalGroups.Network.FibreTyping,
SourceCovers.CausalGroups.Network.CausalGroupTyping,
SourceCovers.CausalGroups.Network.FibreTyping.local_edge_function,
SourceCovers.CausalGroups.Network.FibreTyping.local_output_function}
The computation has a \emph{capacity-conserving causal group typing
conditional on $Z$} if the following holds for each attained $z$.
For each edge $e$ there is one tuple of fields of each type
$a\in[h]$, taking values in an alphabet of size at most
$2^{c_{e,a}(z)}$, where
\[
 c_{e,a}(z)\ge0,\qquad \sum_{a=1}^{h}c_{e,a}(z)\le c_e.
\]
Every outgoing field of type $a$ is a deterministic function only of
incoming type-$a$ fields, the raw bits of $X_a$ in rows attached at that
node, and $z$.  Each output bit in group $a$ is obtained by a function
with these same permitted arguments.  The field alphabets, allocations
and local functions are fixed on the whole fibre $Z=z$, and their
evaluation agrees there with the computation and its outputs.
\end{definition}

The Lean definition expresses local dependence by equality of the
permitted observations on the fibre. Mathlib's function-factorization
theorem supplies actual functions on all local tuples; an attained fibre
supplies default values on unattained tuples. Conversely, local functions
with the stated agreement satisfy this equality-based condition. The
formal network retains the row attachments, ordered outputs and memory
spine. A monotone surjection assigns columns to their nonempty consecutive
groups.

All input-dependent control used by a local function must consequently
be carried by its permitted fields or determined by $z$.  A shared
symbol cannot be reused as a full-capacity field for several groups
without paying the sum of those capacities.  The local functions may be
nonlinear and need only match the original computation on the fibre;
they can be extended arbitrarily to unattained local tuples.  Merely
asserting that an edge's final value depends on one group does not
establish this local causal factorization.

A \emph{partial routing} below uses the underlying undirected network.
For each raw bit it assigns mass in $[0,1]$ to walks from nodes in its
row's attachment set to its own output node.  Every edge traversal,
in either direction and with multiplicity, consumes capacity.  Bits of
unattached rows receive zero mass.  This is an analytic fractional
routing; it is neither a causal execution nor a routing at one common
rate for all bits.

For each group let $\alpha_a$ and $\omega_a$ be its first and last
output nodes, and put
\[
 I_a=\{(v,v+1):\alpha_a\le v<\omega_a\},\qquad
 W_a=db g_a.
\]
The edge intervals $I_a$ are pairwise disjoint, since the groups are
consecutive.  They may share endpoint nodes.  Define the paid volume
\begin{equation}\label{eq:group-volume}
 V=\sum_e c_e+2\sum_{a=1}^{h}W_a(\omega_a-\alpha_a)
 \le \sum_e c_e+2gdb(n-1).
\end{equation}
The factor $2$ pays both collection at $\omega_a$ and return to each
bit's original output node.

\Needspace{12\baselineskip}
\subsection{From autonomous recovery to partial routing}

\begin{theorem}[Theorem Q: paid routing from causal groups]
\label{thm:causal-group-routing}
\formal{SourceCovers/CausalRoutingBound.lean}{SourceCovers.CausalGroups.Network.causal_group_routing,
SourceCovers.CausalGroups.Network.causal_group_routing_sharp}
Assume \cref{def:causal-group-typing}.  For every attained $z$, let
\[
 \mathcal A_a(z)=
 \{x|_{[d]\times J_a\times[b]}:x\in\operatorname{supp}(X\mid Z=z)\},
 \qquad
 R_{\mathrm{sup}}=
 \sum_z\Pr[Z=z]\sum_{a=1}^{h}\log_2|\mathcal A_a(z)|.
\]
There is a partial routing of total mass $R_{\mathrm{sup}}$, with
capacity on edge $e$ at most
\begin{equation}\label{eq:group-augmented-capacity}
 c'_e=c_e+2\sum_{a:e\in I_a}W_a.
\end{equation}
For an integer $\ell\ge1$, let $S_\ell$ be the number of raw bits
whose row has an attachment and whose distance from its nearest row
attachment to its original output node is less than $\ell$.  Then
\begin{equation}\label{eq:causal-group-finite}
 \sum_{a=1}^{h}H(X_a\mid Z)
 \le R_{\mathrm{sup}}
 \le S_\ell+\frac{V}{\ell}.
\end{equation}
In particular, writing $\Delta=L-H(X)$ and
\[
 \operatorname{TC}(X_1;\ldots;X_h\mid Z)
   =\sum_a H(X_a\mid Z)-H(X\mid Z),
\]
one has
\begin{equation}\label{eq:causal-group-helper}
 \Delta+H(Z)\ge
 L+\operatorname{TC}(X_1;\ldots;X_h\mid Z)
   -S_\ell-\frac{V}{\ell}.
\end{equation}
No independence of the input rows, groups, or helper fibres is assumed.
\end{theorem}

\begin{outline}
Represent the input law by weights on its exact support $C$; all these
weights are positive. The helper is indexed by its attained values.
Fix $a,z$.  Causal induction shows that evaluating all fields of type
$a$ depends only on $X_a$ and $z$.  Indeed, this is true of the locally
visible raw data, and it is preserved at each node by the local typing
rule.  Thus a projected input $y\in\mathcal A_a(z)$ defines a code
recovering every bit of $y$, independently of which full input extending
$y$ was used to define it.  This argument evaluates one projection at a
time; it makes no product assumption about the conditional support.

Add a forward collector along each edge of $I_a$. At edge $(v,v+1)$ it
carries precisely the group bits whose original output node is at most
$v$, so it needs only their number of bits in capacity, at most $W_a$.
At node $\omega_a$ this prefix together with the outputs at that node
determines all of $y$. These collector fields still obey group autonomy.
Their exact total capacity is
\[
 \sum_{\xi\in[d]\times J_a\times[b]}(\omega_a-o(\xi))
 \le W_a(\omega_a-\alpha_a),
\]
including zero when both endpoints coincide. The equality follows by
counting, for each output bit, the collector edges that carry it.

Add a supersource $s$ and a source vertex $u_i$ for each row.  Give
$(s,u_i)$ capacity $bg_a$, and give every arc from $u_i$ to a node of
$A_i$ capacity $W_a$.  The sink is $\omega_a$.  Consider any
$s$--$\omega_a$ cut.  If it crosses one of the latter attachment arcs,
its capacity is at least $W_a\ge\log_2|\mathcal A_a(z)|$. Otherwise let
$U$ be its sink-side physical nodes and let $B$ be its sink-side row
vertices.  Every row attached in $U$ belongs to $B$.

For an input $y\in\mathcal A_a(z)$, record its $bg_a$ bits in each
row of $B$ and all type-$a$ fields and collector values entering $U$.
This tuple, with $z$ fixed, determines the computation inside $U$ in
node order.  It supplies every raw argument visible there and every
parent value entering from outside, including any edge that re-enters
$U$ after an earlier departure.  It therefore determines $y$ at the
collector sink.  The tuple map is injective, giving
\[
 \log_2|\mathcal A_a(z)|
 \le bg_a|B|+\text{capacity entering }U.
\]
The right side is the cut capacity. The directed max-flow/min-cut theorem,
whose original undirected form is due to Ford and Fulkerson
\cite{FordFulkerson1956}, supplies a flow of value
$\log_2|\mathcal A_a(z)|$, truncated to that value if necessary.
For completeness, the directed version with real capacities follows by
maximizing flow value over the compact feasible-flow polytope. A residual
walk to the sink gives a signed change on the original edges: divergences
cancel at its internal vertices, and a sufficiently small positive step
respects every capacity bound. Such a walk would increase the maximum.
Without one, the residual reachable set gives a cut of the same value.
Row $i$ supplies at most $bg_a$, and an unattached row supplies zero.

Extract weighted source-to-sink paths carrying the full flow value.
A positive remaining source value forces a path of positive-flow edges,
by conservation across the reachable cut. Subtract the largest permissible
path weight, capped by the remaining source value. This either finishes
or exhausts a positive edge, so at most one path per originally positive
edge is needed. Every traversal is charged with multiplicity.
Remove the paths' artificial source and attachment arcs.
Split each path's weight equally among the $bg_a$ original bits in its row.
Each bit receives that row's total path mass divided by $bg_a$, hence at
most one unit. This preserves every retained edge's load. All these bits
have the same permitted row attachments; no particular bits need to be
identified as entropy carriers.
For each allocated bit append the memory path backwards from
$\omega_a$ to that bit's original output node. This return crosses
$(v,v+1)$ exactly when the bit's output is at most $v<\omega_a$.
The unit bit caps therefore bound the return load by the same prefix
size used by the collector. Their combined charge is at most twice
that prefix size, hence at most $2W_a$ per edge of $I_a$; its total is
at most $2\sum_\xi(\omega_a-o(\xi))$ over the group's bits. Collapsing
collector copies onto original memory edges adds their traversal loads.

Sum over groups.  On an original edge the group flows use at most
$\sum_a c_{e,a}(z)\le c_e$.  The collector and return loads obey
\cref{eq:group-augmented-capacity}.  Average these routings over $z$.
Their common capacity constraints and per-bit mass bounds are linear,
so the average is feasible and has total mass $R_{\mathrm{sup}}$.
Also $H(X_a\mid Z=z)\le\log_2|\mathcal A_a(z)|$, proving the
first inequality in \cref{eq:causal-group-finite}.

At most $S_\ell$ routed mass belongs to short-distance bits, since each
bit has mass at most one.  Every unit of the remaining mass traverses
at least $\ell$ physical edges. Formally, ordinary undirected graph
distance is bounded by each physical walk's length, and total traversal
load equals that length. Taking the infimum over actual row attachments
gives the nearest-attachment distance, with infinity for an unattached row.
Consequently
$\ell(R_{\mathrm{sup}}-S_\ell)\le V$, proving the second inequality.
Unattached rows contribute no routed mass and are not counted as
zero-distance sources.  Finally, because $Z$ is determined by $X$,
\[
 \sum_aH(X_a\mid Z)
 =L-\Delta-H(Z)+\operatorname{TC}(X_1;\ldots;X_h\mid Z).
\]
Substitution proves \cref{eq:causal-group-helper}.
The conditional coordinate-cover inequality also verifies that the total
correlation is nonnegative. The formal theorem retains the sharper prefix
charges above; the stated interval-budget version follows directly.
Finally, disjoint collector links inject into the $n-1$ memory slots,
giving the last inequality in \cref{eq:group-volume}.
\end{outline}

The theorem supplies an explicit conversion from a restricted causal
code to partial routing.  Its extra communication is fully charged.
It does not strengthen the general routing obstruction, and it does
not convert arbitrary causal codes to routings.

\subsection{Application to linear-length tape diagrams}

For clarity we state the precise short-distance estimate needed for
partial, possibly unequal, commodity masses.  The summed-distance bound
alone would not suffice.

\Needspace{10\baselineskip}
\begin{lemma}[Attached short-distance count]\label{lem:attached-short-count}
\formal{SourceCovers/NativeShortCount.lean}{SourceCovers.attachedShortBits, SourceCovers.card_attachedShortBits_le_local, SourceCovers.card_attachedShortBits_le_profile, SourceCovers.card_attachedShortBits_le, SourceCovers.attachedShortBits_zero_radius, SourceCovers.attachedShortBits_zero_block}
Let a $k$-work-tape diagram have block length $r\ge0$, $n$ block
nodes, and raw length $L=d^2b$. Use its original row attachments and
transposed-bit output nodes. Write $R(a)$ for the rows attached at node
$a$, $B_t(a)$ for its closed undirected radius-$t$ ball, and $m_{i,s}$
for the number of bits of raw row $i$ assigned to output node $s$.
For every integer $\ell\ge1$,
\begin{equation}\label{eq:attached-short-local}
 S_\ell\le
 \sum_{a<n}\sum_{i\in R(a)}\sum_{s\in B_{\ell-1}(a)}m_{i,s}.
\end{equation}
This bound requires neither diagram validity nor an upper bound on $r$.
If the diagram is valid for output length $L$ and $r\le db$, then
\begin{align}
 S_\ell&\le b\sum_{a<n}|R(a)|\,|B_{\ell-1}(a)|,
     \label{eq:attached-short-profile}\\
 S_\ell&\le\min\{L,\,(2n-1)_+b(2k+3)^{\ell-1}\}.
     \label{eq:attached-short-count}
\end{align}
Also $S_0=0$, and $S_\ell=0$ for $r=0$ under the library's zero-block
convention. The Lean inequalities include these cases, interpreting
natural subtraction $\ell-1$ as $\max\{\ell-1,0\}$ and
$(2n-1)_+=\max\{2n-1,0\}$.
\end{lemma}

\begin{outline}
The formal definition excludes unattached rows, whose inherited default
distance is zero. If an attachment exists, the block network is nonempty,
its chosen output node lies in that network, and the memory spine
connects it to the attachment. These facts require no validity
hypothesis. The bridge lemma therefore identifies a short bit with a
row attachment $a$ whose radius-$(\ell-1)$ ball contains its output
node. Taking the union over $a$, its attached rows, and its ball gives
\cref{eq:attached-short-local} by Mathlib's finite-union bound.

For a valid diagram and $1\le r\le db$, a block's output positions
lie in an interval of length at most $db$. Two distinct bits of one raw
row with the same within-entry offset have transposed positions separated
by a nonzero multiple of $db$. They cannot both lie in this interval.
The within-entry offset thus injects the bits counted by $m_{i,s}$ into
$\{0,\ldots,b-1\}$, proving $m_{i,s}\le b$ and
\cref{eq:attached-short-profile}. This sharpens the first paper's
$2b$ estimate, which separately counts the portions in two output
columns~\cite[section ``Codes on time--space diagrams'']{FirstPaper2026}.

A block attaches at most two raw rows, and its underlying graph has
degree at most $2k+2$. Reuse the first paper's
\lean{card_rowsAt_le} and \lean{card_ball_le} from
\leanfile{Transpose/Tape/CECount.lean} to obtain
$|R(a)|\le2$ and $|B_t(a)|\le(2k+3)^t$.
The first block can touch only row zero: the input head starts at zero
and has fewer than $r\le db$ steps in that block. Consequently
$\sum_{a<n}|R(a)|\le(2n-1)_+$. Together with $S_\ell\le L$,
this gives \cref{eq:attached-short-count}. The zero-radius set is empty by its
strict distance test; for zero block length there are no nodes or
attachments.
\end{outline}

The profile bound retains actual row and ball counts for later use,
without imposing a common routing rate on the counted bits. The sharper
coefficient and first-block correction improve the finite estimate;
they do not change the asymptotic obstruction below. No optimality of the
graph-ball estimate or the union bound is claimed.

\Needspace{8\baselineskip}
\begin{corollary}[Narrow causal groups require an almost complete helper]
\label{cor:causal-group-obstruction}
\notformalized
Fix $k\ge0$, an integer $G\ge1$, and $c>0$.  Along any sequence
$d\to\infty$, let $b\ge1$, $L=d^2b$, and
$r=\lceil db/G\rceil$.  Suppose valid $k$-work-tape diagrams have
length $T\le cL$, and their code capacities satisfy the block-profile
volume bound
\[
 \sum_e c_e\le4knrG,
 \qquad n=\lceil T/r\rceil.
\]
Let $X$ be any law on correctly transposed inputs and $Z$ a
raw-input-determined helper admitting causal group typing into fixed
consecutive column groups of maximum size $g$.  Then, with constants
depending only on $k,G,c$,
\begin{equation}\label{eq:causal-group-asymptotic}
 \sum_aH(X_a\mid Z)
 =O\!\left(\frac{L}{\sqrt d}+\frac{(1+g)L}{\log d}\right).
\end{equation}
In particular, if $g=o(\log d)$, then
\[
 \Delta+H(Z)\ge L-o(L).
\]
Thus $\Delta=o(L)$ and $H(Z)=o(L)$ cannot both hold.  The conclusion
applies in particular to the series' standard width
$b=b(d)=\max\{1,\lceil\log_2d\rceil\}$.
\end{corollary}

\begin{informalproof}
Since $r\ge db/G$, one has $n\le cGd+1=O(d)$.  Also
$nr\le T+r$ and $rG\le db+G$, so
$\sum_e c_e=O(L)$.  Equation~\eqref{eq:group-volume} therefore gives
$V=O((1+g)L)$.
For sufficiently large $d$, take
$\ell=\lfloor\tfrac12\log_{2k+3}d\rfloor\ge1$.
Then $\ell$ is of order $\log d$, and
\cref{lem:attached-short-count} gives
$S_\ell=O(db\sqrt d)=O(L/\sqrt d)$.
Apply \cref{thm:causal-group-routing} to obtain
\cref{eq:causal-group-asymptotic}.  Finally,
$L-\Delta-H(Z)=H(X\mid Z)\le\sum_aH(X_a\mid Z)$.
The right side is $o(L)$ when $g=o(\log d)$.
\end{informalproof}

At the standard width, the block-profile volume hypothesis is supplied
by the Lean-verified \lean{sum_cap_le} in~\cite{FirstPaper2026}.
The finite theorem also
permits other profiles: a separate $O(L)$ total capacity allowance
leaves the asymptotic conclusion unchanged.

For the construction program, the corollary excludes a completion that
preserves one fixed narrow consecutive grouping at every local update
after a sublinear-entropy joint helper is fixed.  Arbitrarily many
updates and arbitrary nonlinear operations inside each group remain
allowed.  It does not exclude final column purity, which every correct
transposer already has, nor a succession of small groups that migrate
and cumulatively mix all columns.  For nonconsecutive groups the exact
collector cost in \cref{eq:group-volume} is still available, but its
$2gdb(n-1)$ upper bound need not hold.  Finally, none of the geometric
hypotheses elsewhere in this paper is asserted to produce this causal
typing with a cheap helper; $Z=X$ always gives a vacuous expensive one.
